\section{Methodology}
\label{sec:method}

We formalize the online routing optimization of Section~\ref{sec:problem:objective} as an MDP (\S\ref{sec:method:mdp}), justify the choice of PPO as the training algorithm (\S\ref{sec:method:algo}), and present the flow-aware QoS-FiLM architecture (\S\ref{sec:method:overview}--\ref{sec:method:train}). All network dimensions, layer counts, and training coefficients are written symbolically here; their specific values are listed in Section~\ref{sec:exp:hyperparam}.

\subsection{Architecture Overview}
\label{sec:method:overview}

Fig.~\ref{fig:film_gnn_arch} gives a high-level picture of the QoS-FiLM architecture, which consists of three main components. \textbf{(1) FiLM-GNN encoder:} takes as input the current graph snapshot (node and link features) and the flow intent vector $m_t = [\phi_t, \tilde b_t]$. At each graph attention layer, the intent vector generates per-channel scaling $\gamma^{(\ell)}$ and shifting $\beta^{(\ell)}$ parameters via lightweight MLPs; these modulate the message-passing output through an affine transformation $x^{(\ell)} = (1+\gamma^{(\ell)})\odot \tilde x^{(\ell)} + \beta^{(\ell)} + x^{(\ell-1)}$, so that the same physical topology yields different node embeddings for flows with distinct QoS requirements. \textbf{(2) Actor Head:} pools the node embeddings along each candidate path via an LSTM-based PathPooler, concatenates the resulting path representation with $m_t$, and scores the $K$ candidates through the KPathScorer to produce a Categorical action distribution. \textbf{(3) Critic Head:} aggregates the node embeddings into a graph-level vector via attention pooling (where attention weights depend on $m_t$) and regresses the state value $V_\psi(s_t)$ through an MLP. These three components operate in a closed MDP loop: the agent observes the current state, selects a path for the arriving flow, the environment updates link capacities, and PPO optimizes the policy from the collected trajectories.

\begin{figure}[t]
\centering
\includegraphics[width=\linewidth]{ENV-REACT.pdf}
\caption{Overall architecture of QoS-FiLM.}
\label{fig:film_gnn_arch}
\end{figure}

\subsection{Problem Formalization: MDP Modeling}
\label{sec:method:mdp}

We formalize the online routing optimization of Section~\ref{sec:problem:objective} as a Markov decision process $(\mathcal{S},\mathcal{A},P,r)$. The agent makes sequential decisions: at step $t$ it selects, for the $t$-th flow $\mathbf{f}_t$, a path from the candidate set $\mathcal{P}_t$.

\textbf{State} $s_t = (\mathcal{G}_t,\, m_t,\, \mathcal{P}_t) \in \mathcal{S}$ consists of three parts:
\begin{enumerate}
\item \emph{Graph snapshot} $\mathcal{G}_t$: node features (marking the source and destination $o_t,v_t$ of the current flow) and link features $[d_e,\,c_e,\,b_e^{\text{res}}]$ (delay, capacity, residual bandwidth);
\item \emph{Flow intent vector} $m_t=[\phi_t,\,\tilde b_t]$, where $\tilde b_t$ is the normalized bandwidth demand;
\item \emph{Candidate paths} $\mathcal{P}_t$ and their validity masks: each path is padded to a fixed tensor length, and the binary mask indicates which positions correspond to actual path nodes versus padding tokens.
\end{enumerate}
The observation contains only the current flow $\mathbf{f}_t$ and excludes future flows $\{\mathbf{f}_\tau\}_{\tau>t}$, consistent with the online causal constraint of Section~\ref{sec:problem:objective}.

\textbf{Action} $a_t\in\mathcal{A}=\{0,1,\dots,K-1\}$ is a discrete choice, selecting the $a_t$-th candidate path $p_t=p_t^{a_t}$ from the $K$ candidates to carry the current flow.

\textbf{Transition} must be distinguished at two levels. Given the current flow and the selected path, \emph{resource accounting is deterministic and known}: the residual bandwidth $b_e^{\text{res}}$ of the links along the path is reduced by $\hat b_t$ (Eq.~\eqref{eq:rho}, the actually allocated bandwidth) and the pointer advances. The next arriving flow $\mathbf{f}_{t+1}$, however, is drawn from the demand distribution $\mathcal{D}$, is invisible to the agent before it arrives, and is unaffected by the current action. The full state transition $P(s_{t+1}\mid s_t,a_t)$ is therefore not deterministic (given $s_t,a_t$, the next state $s_{t+1}$ is not unique); the stochasticity comes from the arrival of the next flow. This matches the ``deterministic accounting, unseen future'' structure described in Section~\ref{sec:problem:objective}.

\textbf{Reward} is taken directly from the utility function $U(\mathbf{f}_t,p_t)$ of Section~\ref{sec:problem:objective} (Eq.~\eqref{eq:utility}), shifted by a constant $-1$ to narrow the reward range and facilitate value function learning:
\begin{equation}
r_t = U(\mathbf{f}_t, p_t) - 1
    = -1 + U^{\text{bw}}_t + U^{\text{dl}}_t.
\label{eq:reward}
\end{equation}
The reward range is $[-1, 1]$. For elephant flows ($\phi_t\to 1$) the delay utility term approaches zero and the reward upper bound is $0$ (attained when $\rho_t=1$); for mouse flows ($\phi_t\to 0$), when bandwidth is fully satisfied and delay is very low, both utility terms can take relatively large positive values, so the reward for this class of flows can exceed $0$. Under a test distribution dominated by mouse flows, the mean reward can therefore be positive overall. Because the offset constant does not change the optimal policy, maximizing the cumulative reward $\sum_t r_t$ is equivalent to maximizing the cumulative utility (Eq.~\eqref{eq:objective}). The structure of Eq.~\eqref{eq:reward} implies that, to earn high rewards on both mouse flows (which require hard bandwidth satisfaction and low latency) and elephant flows (which value the throughput ratio), the policy must \emph{perceive and differentiate the QoS intent} $\phi_t$ of each flow. This motivates the flow-aware encoder presented below.

\subsection{Choice of Training Algorithm}
\label{sec:method:algo}

Section~\ref{sec:problem:objective} argued that this problem is an online stochastic optimization over the demand distribution that requires a non-myopic policy. We adopt PPO as the training algorithm for several reasons: its stochastic policy naturally parameterizes the discrete choice among $K$ candidate paths through a Categorical distribution; its clipped objective keeps updates stable (accommodating the large difference in reward scales between the two flow types); on-policy data collection matches the online deployment setting; combined with generalized advantage estimation (GAE) it achieves the long-range credit assignment that non-myopic decision-making demands; and, being model-free, it does not require explicit modeling of the transition stochasticity caused by the unknown arrival distribution $\mathcal{D}$.

\subsection{Flow-Aware QoS-FiLM Encoder}
\label{sec:method:encoder}

The encoder maps the graph snapshot in the state into node embeddings (i.e., a vector representation for each node). The core idea is that every layer of message passing should ``know which flow is currently being routed,'' so that flows with different QoS intents yield different topological representations. We implement this conditioning with FiLM~\cite{Perez2018_FiLM}. Let the link feature vector be $f_e = [d_e,\, c_e,\, b_e^{\text{res}}] \in \mathbb{R}^3$, the concatenation of delay, capacity, and residual bandwidth.



Let the hidden dimension be $h$. An MLP $\mathrm{MLP}_{\text{lift}}$ first lifts the raw node features to $h$ dimensions: $x^{(0)}=\mathrm{MLP}_{\text{lift}}(x)$. We then stack $L$ graph attention blocks.

The $\ell$-th layer receives the current flow's intent vector $m_t$ and generates per-channel (per-dimension of the $h$-dimensional hidden vector) scaling and shifting parameters through two independent MLPs:
\begin{align}
\gamma^{(\ell)} &= \mathrm{MLP}^{(\ell)}_{\gamma}(m_t), \quad \gamma^{(\ell)}\in\mathbb{R}^{h},
\label{eq:film_gamma}\\
\beta^{(\ell)} &= \mathrm{MLP}^{(\ell)}_{\beta}(m_t), \quad \beta^{(\ell)}\in\mathbb{R}^{h},
\label{eq:film_beta}
\end{align}
and applies FiLM modulation to the message-passing output:
\begin{align}
\tilde x^{(\ell)} &= \mathrm{ELU}\!\big(\mathrm{LN}(\mathrm{Conv}^{(\ell)}(x^{(\ell-1)}, E, f))\big),
\label{eq:conv}\\
x^{(\ell)} &= \big(1+\gamma^{(\ell)}\big)\odot \tilde x^{(\ell)} + \beta^{(\ell)} \;+\; x^{(\ell-1)},
\label{eq:film}
\end{align}
where $\mathrm{Conv}^{(\ell)}$ is a graph attention operator (default: TransformerConv~\cite{Shi2021_UniMP}
% TransformerConv 引用查证:PyTorch Geometric 官方文档及源码(transformer_conv.py)
% 均标注该算子来自 "Masked Label Prediction: Unified Message Passing Model for
% Semi-Supervised Classification"(Shi et al., IJCAI 2021, arXiv:2009.03509),
% 即 Shi2021_UniMP。经三方交叉验证确认此为原始文献。
-style dot-product attention,
with link features $f$ as edge attributes), $\mathrm{LN}$ is layer normalization, $\odot$ is element-wise multiplication, and the final term is a residual connection (adding the layer input directly to its output to ease the training of deep networks). Eq.~\eqref{eq:film} is the key to our method: at \emph{every layer}, the intent vector $m_t$ reshapes node representations through the affine transformation $(1+\gamma)\odot\tilde x+\beta$, so that the same physical topology is encoded into different embeddings for mouse flows and elephant flows. The fundamental distinction from GNN-FiLM~\cite{Brockschmidt2020_GNNFiLM} is the source of $\gamma,\beta$: GNN-FiLM computes them from the \emph{target node representation} to modulate aggregation, whereas we compute them from the \emph{per-flow QoS intent}, making message passing flow-aware.

\subsection{Path Pooling, Scoring, and Value Network}
\label{sec:method:scorer}

\textbf{Path Pooling.} A path pooler (PathPooler) converts the node embeddings $x^{(L)}$ into candidate path representations. For the $k$-th candidate path, the embeddings along its node sequence are extracted to form a variable-length sequence, encoded by an LSTM whose final hidden state is taken, and then passed through layer normalization to yield the path embedding $z_k\in\mathbb{R}^{h}$. The $K$ candidate paths are thus pooled into $\{z_k\}_{k=1}^{K}$.

\textbf{Scoring and Policy.} The scorer $\mathrm{KPathScorer}$ concatenates each path embedding with the current flow intent as $[z_k; m_t]$, passes it through a feature MLP (ELU $+$ LayerNorm) and then a scoring-head MLP, and outputs a scalar logit $\ell_k$ (an unnormalized score). The $K$ logits form a Categorical distribution from which the policy samples an action:
\begin{equation}
\pi_\theta(a_t\!=\!k\mid s_t) = \mathrm{softmax}_k(\ell_1,\dots,\ell_K).
\label{eq:policy}
\end{equation}
If the intent vector $m_t$ entered only at the final scorer, the paths would already have been pooled into fixed vectors, so intent information could no longer back-propagate to influence the topological representation, causing irreversible information loss.

\textbf{Value Network.} The policy and value networks use independent encoders because their optimization objectives conflict: the policy must learn relative preferences among actions to support exploration, whereas the value function must regress scalar returns accurately. Sharing parameters would cause gradient interference and training instability. The value network pools its output node embeddings into a single graph-level vector via attention: each node's attention weight is computed from its embedding and the flow intent $m_t$, so that nodes relevant to the current flow's QoS requirements contribute more to the pooled representation; an MLP regresses this vector into a scalar state value $V_\psi(s_t)$.

\subsection{Training Procedure}
\label{sec:method:train}

Training uses PPO, with advantages computed by GAE. The loss combines the clipped policy objective $\mathcal{L}^{\text{clip}}_\theta$, the value loss $c_v\mathcal{L}^{\text{value}}_\psi$, and the negative entropy bonus $-c_e\mathcal{H}[\pi_\theta]$:
\begin{equation}
\mathcal{L}(\theta,\psi) = \mathcal{L}^{\text{clip}}_\theta + c_v\,\mathcal{L}^{\text{value}}_\psi - c_e\,\mathcal{H}[\pi_\theta],
\label{eq:loss}
\end{equation}
where $\mathcal{L}^{\text{clip}}_\theta$ is the policy objective controlled by the clipping coefficient $\epsilon$, $c_v, c_e$ are value and entropy coefficients respectively, and advantages are normalized within each minibatch. The complete procedure is given in Algorithm~\ref{alg:ppo}.

\begin{algorithm}[t]
\caption{PPO Training for QoS-FiLM}
\label{alg:ppo}
\SetKwInOut{Input}{Input}
\SetKwInOut{Output}{Output}
\Input{Network $G$; demand distribution $\mathcal{D}$; candidate count $K$; sequence length $W$; total frames $T$; batch frames $B$; PPO epochs $E$; learning rate $\eta$}
\Output{Policy parameters $\theta$, value parameters $\psi$}
Randomly initialize $\theta,\psi$\;
\While{collected frames $<T$}{
  \tcp{Collection (on-policy)}
  Interact with environment using current policy $\pi_\theta$ to collect $B$ frames; per step: encode graph $\to$ pool $K$ paths $\to$ score $\to$ sample action via $\pi_\theta$ $\to$ execute routing, deduct bandwidth, record reward $r_t$ (Eq.~\eqref{eq:reward})\;
  \tcp{Advantage estimation}
  Compute advantages $\hat A_t$ and return targets $\hat R_t$ using $V_\psi$ and GAE\;
  \tcp{Optimization}
  \For{$E$ PPO epochs}{
    \ForEach{minibatch}{
      Normalize $\hat A_t$ within minibatch\;
      Compute loss $\mathcal{L}(\theta,\psi)$ (Eq.~\eqref{eq:loss})\;
      Clip gradients, update $\theta,\psi$ with Adam (learning rate $\eta$)\;
    }
  }
  Update learning rate via cosine annealing\;
  \tcp{Evaluation and checkpointing}
  Run several rollouts with deterministic policy, record average reward; save best checkpoint\;
}
\end{algorithm}
